\unitheading{K10\quad 流体运动}{unit-k10}
\addcontentsline{toc}{chapter}{K10 流体运动}
\label{ch13:sec:flow-intro}

水和空气没有固定形状，受力后可以流动，统称流体。流体运动极其复杂，本节先讨论稳定流动：空间中任一给定点的速度大小和方向都不随时间改变。稳定并不表示某个流体微元的速度沿途不变；微元运动到不同位置时，会取得当地不同的流速。

流体由大量分子组成，逐一跟踪每个分子的运动通常既不可行也没有必要。稳定流动改为考察“空间中某一点此刻允许多大的速度、指向哪里”：只要这个速度场不随时间变化，后来经过同一点的流体微元就会取得相同速度。某个微元沿途经过不同位置时仍可能不断加速、减速或转弯，所以“稳定”描述的是空间分布不随时间变，并不等于每个微元作匀速直线运动。

水在通常条件下密度变化很小，可理想化为不可压缩流体；空气的密度较易随压强和速度改变，是可压缩流体。对稳定流动，给出空间每点的速度矢量便给出流速分布。沿速度场逐步跟踪微元可得到流线。

不可压缩同样是一种近似：它表示同一流体微元在运动中密度近似不变，并不是说流体“不能改变形状”。气体通常更容易压缩，完整描述稳定气流时除了速度场，还须给出各点密度。后面的连续性方程会分别说明体积流率和质量流率在两种模型中的区别。

\SphereStreamlinesDiagram

在速度场连续光滑且速度非零的区域，流线具有三条基本性质：流线在每一点与当地速度相切；每一点通过一条局部唯一的流线；两条不同流线不能相交。停滞点的速度为零，不能由零矢量指定切线方向，需作为例外单独讨论。

流线可以用逐步推进的方法构造：从一个起点出发，沿当地速度方向前进很短距离，到达新位置后改沿新位置的速度方向继续前进；令每一步趋于无限小，得到的光滑曲线就是流线。图2.23中的每条线都应在圆球附近平滑转弯而不穿过球体。若两条流线相交，交点处会出现两个切线方向，与稳定速度场在每点只有一个速度矢量矛盾。

取一个与流动方向垂直的面积 $S$，若截面上流速均为 $v$，单位时间通过的体积为 $Sv$。由一圈流线围成的管状区域称为流线管，其侧壁没有流体穿过。

在短时间 $\Delta t$ 内，穿过截面的流体近似占据底面积 $S$、长度 $v\Delta t$ 的柱体，体积为 $Sv\Delta t$，因此体积流率为 $Sv$。流线管的侧壁处处由流线组成，速度与侧壁相切，没有法向分量，所以流体不会从侧壁进出；流入和流出只能通过两端截面。

\StreamtubeDiagram

图~\ref{ch13:fig:k10-2-24}的两截面属于同一流线管。不可压缩流体稳定流过同一流线管时，单位时间通过各截面的体积相等。截面积增大则流速降低，截面积减小则流速增大。

\begin{origexamplebox}{例2.97\quad 不可压缩流体的连续性方程}
不可压缩流体在管中稳定流动，每个垂直截面上流速均匀。截面 $A$ 的面积为 $S_A$、流速为 $v_A$；截面 $B$ 的面积为 $S_B$。求 $v_B$。

\PipeContinuityDiagram
\end{origexamplebox}
\begin{xesolution}
图~\ref{ch13:fig:k10-ex-2-97}标出截面 $S_A,S_B$ 与对应速度。取相同时间 $\Delta t$，通过两截面的体积分别为 $S_Av_A\Delta t$、$S_Bv_B\Delta t$。稳态且不可压缩时两者相等：
\[
S_Av_A=S_Bv_B.
\]
所以
\[
\boxed{v_B=\frac{S_A}{S_B}v_A}.
\]
量纲检查：面积比无量纲，结果仍为速度。
\end{xesolution}

\begin{origexamplebox}{例2.98\quad 等截面管内能否稳定加速}
圆柱形管道处处充满不可压缩流体。流体能否在管中作稳定的加速运动？采用沿直管轴线流动、各横截面流速均匀、无支流的一维模型。

\UniformPipeDiagram
\end{origexamplebox}
\begin{xesolution}
图~\ref{ch13:fig:k10-ex-2-98}中圆柱管的截面积 $S$ 处处相同。连续性方程要求任意两截面
\[
Sv_1=Sv_2,
\]
故 $v_1=v_2$。在“不可压缩、稳定、处处充满且无支流”的条件下，流速不能沿管长增加，因而
\[
\boxed{\text{不能作沿程稳定加速运动}}.
\]
若流体可压缩、管道有泄漏或流动随时间变化，结论可能改变；这些都超出了题设。
若不作一维近似，连续性只保证截面平均流速相同；截面内部仍可发生速度重新分布。因此不能由“等截面”一句话推出任何微元的局部加速度都为零。
\end{xesolution}

\begin{origexamplebox}{例2.99\quad 下落水束的截面积}
极大容器底部有面积为 $S$ 的圆孔，水刚离孔时速度为 $v$。每个水分子随后自由落体，水束不散开。求孔下方 $h$ 处水流截面积 $S(h)$。

\FallingJetDiagram
\end{origexamplebox}
\begin{xesolution}
图~\ref{ch13:fig:k10-ex-2-99}中，$h$ 的下端、$S(h)$ 与 $v(h)$ 标在同一横截面。自由下落的速度由
\[
v(h)^2=v^2+2gh
\]
给出，故 $v(h)=\sqrt{v^2+2gh}$。水不可压缩，体积流率不变：
\[
Sv=S(h)v(h).
\]
因此
\[
\boxed{S(h)=\frac{Sv}{\sqrt{v^2+2gh}}}.
\]
随 $h$ 增大，水束加速而截面收细，乘积 $S(h)v(h)$ 保持为 $Sv$。
\end{xesolution}

对于可压缩流体，稳定流动中真正守恒的是质量流率而非体积流率。若没有物质产生或消失，对任意闭合曲面，单位时间流入的质量必须等于流出的质量，否则曲面内的质量会随时间积累或减少。这就是连续性条件。

若截面上密度和速度近似均匀，短时间通过的质量为 $\rho Sv\Delta t$，质量流率为 $\rho Sv$。不可压缩情形中同一流线管内 $\rho$ 不变，才可约去密度得到 $Sv$ 恒定；可压缩流体则可能用密度变化补偿截面积或速度变化。连续性只表达物质守恒，不自动给出压强和速度的全部关系，解决更完整的流体问题还需要后续动力学方程。

连续性方程的适用对象必须选成同一股稳定流。若管道存在支路，应把所有入口和出口的质量流率带符号相加，而不能只随意比较两个截面；若流动随时间变化，控制体内还可能发生质量积累。对题图的单一流线管，侧壁无穿流、内部又无源汇，所以任意两个横截面的流率才相等。

自由水束的两个例题展示了同一守恒条件的两种实现方式。真实水近似不可压缩，密度保持不变，下落加速后必须缩小截面积；若人为规定截面积不变而又让各微元自由加速，连续性就要求密度降低。二者不能同时把截面积和密度都固定，否则 $v$ 增大将使质量流率凭空增大，违反稳定流动的物质守恒。

定常流动中，流线与单个微元的轨迹重合；非定常流动中一般不重合。若速度场随时间改变，某一时刻画出的流线仍处处与当时速度相切，但一个微元随后看到的速度场已经变化，它实际走出的迹线可能与这条瞬时流线不同。本节始终假定稳定流动，因此可以用流线同时表示当地速度方向和微元运动轨迹；超出这一条件时应区分两者。

截面若不垂直于速度，穿过它的体积流率应使用法向分量 $v_n$，即 $Sv_n$，而不是总速率 $Sv$。若截面上速度或密度不均匀，还需把 $\rho\vect v\cdot\vect n$ 在整个截面上积分。本节例题把截面选为垂直于流动且各处均匀，才得到简洁的 $Sv$ 或 $\rho Sv$；这些条件也是公式的一部分，不能从图中省略。

\begin{origexamplebox}{例2.100\quad 可压缩流体的质量连续性}
流体在管道中稳定流动，且同一垂直截面上性质均匀。$A$ 处截面积、密度、流速分别为 $S_A,\rho_A,v_A$；$B$ 处为 $S_B,\rho_B,v_B$。求 $v_B$。

\CompressiblePipeDiagram
\end{origexamplebox}
\begin{xesolution}
图~\ref{ch13:fig:k10-ex-2-100}将两侧密度分别标为 $\rho_A,\rho_B$。单位时间通过截面的质量为 $\rho Sv$，不能只比较体积。质量连续性给出
\[
\rho_AS_Av_A=\rho_BS_Bv_B,
\]
所以
\[
\boxed{v_B=\frac{\rho_AS_A}{\rho_BS_B}v_A}.
\]
当 $\rho_A=\rho_B$ 时退化为例2.97的不可压缩结果。
\end{xesolution}

\begin{origexamplebox}{例2.101\quad 等截面自由落体液柱的密度}
极大容器中液体初始密度为 $\rho$，从面积为 $S$ 的孔以速度 $v$ 流出。忽略分子间相互作用和表面张力，设下落过程中水束截面积始终为 $S$，每个分子作自由落体。求密度 $\rho(h)$。
\end{origexamplebox}
\begin{xesolution}
下落 $h$ 后速度仍为
\[
v(h)=\sqrt{v^2+2gh}.
\]
稳态质量流率在孔口和高度 $h$ 处相等：
\[
\rho Sv=\rho(h)S v(h).
\]
截面积约去，得到
\[
\boxed{\rho(h)=\frac{\rho v}{\sqrt{v^2+2gh}}}.
\]
这里人为固定了截面积，因此加速必须由密度降低来满足质量连续性；这正与例2.99中“密度不变、截面积减小”互为对照。
\end{xesolution}
